\scanpage{12}
\mspage{1)}
Espaces stratifiés (étude par voisinages \deleted{tubulaires}
\inserted{coniques})
\textcircled{I}

\textcircled{1} Soit $Y$ un espace topologique. Un \uline{voisinage \deleted{tubulaire}} \inserted{\uline{\unsure{conique}}}
\inserted{(``\uline{absolu}'')} de \deleted{$Y$}
\inserted{(ou ``cône sur $Y$'')} est un espace topologique
$T\deleted{\unsure{_Y}}$ contenant $Y$,
(le ``\uline{tube}''), muni d'une partie fermée $\dot T$
\deleted{\illegible{1}} (le ``bord'' du tube'') et d'une opération du monoïde topologique
$\mathbf{I}=[0,1]\,\unsure{}$ (pour la multiplication) sur $T$, notée
% P08: present overworked mark after [0,1]; identity and cancellation state uncertain.
\[
(1)\qquad (\lambda,t)\longmapsto\lambda t
\qquad \mathbf{I}\times T\longrightarrow T
\]
satisfaisant :
\[
\left.\begin{array}{ll}
a)&\lambda y=y\quad\forall y\in Y\\
b)&0.\deleted{\unsure{T|\bar T}}T\subset Y
\end{array}\right\}
\]
\margin{donc les $\lambda_T:T\longrightarrow T$
$(0\leq\lambda\leq1)$ forment une homotopie de
l'identité $\operatorname{id}_T$ (pour $\lambda=1$)
sur une \uline{rétraction} de $T$ sur $Y$,
laissant les pts de $Y$ fixes au cours de l'homotopie}
\margin{\uline{NB} La condition b) sous \unsure{sa|la} forme faible
$0.\dot T\subset Y$, plus la condition c)
(à qui celle-ci donne un sens) suffisent :
entraînent}

c) Le \deleted{application} diagramme d'applications
\diagram{\[
\unsure{(2)|(b)}\qquad
\begin{tikzcd}[column sep=large,
execute at end picture={
\path (p012a.east) -- coordinate[pos=.32] (p012extrahead) (p012b.west);
\draw ([xshift=-3pt,yshift=2pt]p012extrahead) -- (p012extrahead) -- ([xshift=-3pt,yshift=-2pt]p012extrahead);
}]
|[alias=p012a]| \mathbf{I}\times\dot T \arrow[r,"{(\lambda,\unsure{x})\mapsto\lambda x}",
"\deleted{\unsure{g|f}}"'] & |[alias=p012b]| T\\
\{0\}\times\dot T \arrow[u,<->] \arrow[r,"{(0,x)\mapsto0.x}"'] & Y \arrow[u,hook,"\mathrm{inc}"']
\end{tikzcd}
\]}
est \uline{cocartésien}$^{*}$, i.e. identifie $T$ au cône de l'application
\[
(3)\qquad f:\dot T\longrightarrow Y\qquad f(\unsure{a|x})=0.x
\]
\margin{\uline{$^{*}$} Cela signifie que
$\mathbf{I}_{]0,1]}\times\dot T\xrightarrow{\substack{\mathrm{bij.}\\\sim}}T\setminus Y$
et qu'une partie de $T$ est ouverte ssi
\deleted{\unsure{son intersection}} \inserted{\unsure{ses images}}
\unsure{et} \unsure{dans $Y$ et ses images dans $\mathbf{I}\times\dot T$}}
\[
(4)\qquad T\simeq C(f)
\]
Dans cette identification, $Y\hookrightarrow\deleted{\unsure{T}}T
\hookleftarrow\dot T$ correspondent
aux applications canoniques \deleted{\unsure{du cône $f$}}
$\operatorname{but}(f)\hookrightarrow C(f)\hookleftarrow
\operatorname{source}\unsure{(}f\unsure{)}$
et l'opération de $\mathbf{I}$ sur $T$ provient de l'opération de
$\mathbf{I}$ sur $C(f)$, définie par passage au quotient de
l'opération évidente de $\mathbf{I}$ sur $\mathbf{I}\times\dot T$, et l'op. triviale
de $\mathbf{I}$ sur $Y$. Les axiomes a) b) c) signifient
aussi que la situation $(Y\mathrel{\tikz[baseline=-.45ex]{\draw (.12,.09) .. controls (-.01,.09) and (-.01,0) .. (.12,0) -- (.65,0);}}T
\mathrel{\tikz[baseline=-.45ex]{\draw (0,0) -- (.53,0) .. controls (.68,0) and (.68,.11) .. (.53,.13);}}\dot T,
\mathbf{I}\times T-T)$
est isomorphe : celle d'un $C(f)$ avec $f:\dot T-Y$.
Notons que $\dot T$ est connu quand on connaît l'opération
de $\mathbf{I}$ sur $T$, \deleted{\unsure{comme on}} en posant
\[
(5)\qquad \mathring T=T\setminus\dot T
\quad\text{(l'\uline{intérieur} du tube)}
\qquad\text{et notons que }\mathring T=\bigcup_{\lambda\in[0,1[}\lambda T
\]
\margin{11}
